\subsection{同调维数的局部化}

命题 2.--- 设 A 为环，M 与 N 为 A-模，i 为整数，S 为 \(A\) 的一个乘性子集。我们有 \(\mathrm{S}^{-1}\mathrm{A}\) -模的一个典范同构

\[
\mathrm{S} ^ {- 1} \operatorname{Tor} _ {i} ^ {A} (\mathrm{M}, \mathrm{N}) \longrightarrow \operatorname{Tor} _ {i} ^ {\mathrm{S} ^ {- 1} A} (\mathrm{S} ^ {- 1} \mathrm{M}, \mathrm{S} ^ {- 1} \mathrm{N}).
\]

若环 A 是 Noether 的且 A-模 M 有限生成，则有 S \(^{-1}\) A-模的一个典范同构

\[
\mathrm{S} ^ {- 1} \operatorname{Ext} _ {A} ^ {i} (\mathrm{M}, \mathrm{N}) \longrightarrow \operatorname{Ext} _ {\mathrm{S} ^ {- 1} \mathrm{A}} ^ {i} (\mathrm{S} ^ {- 1} \mathrm{M}, \mathrm{S} ^ {- 1} \mathrm{N}).
\]

由于 A-模 \(S^{-1}A\) 是平坦的，这由 A, X, 第 110 页, 命题 9 及第 111 页, 命题 10 得出。

推论.--- 设 A 为环，M 与 N 为 A-模，i 为整数。

\begin{enumerate}
\def\labelenumi{\alph{enumi})}
\item
  \(\mathrm{Tor}_{i}^{A}(\mathrm{M},\mathrm{N})\) 的支撑集含于 \(\mathrm{Supp}(\mathrm{M})\cap \mathrm{Supp}(\mathrm{N})\)，且当 A 是 Noether 的而 M 有限生成时，\(\mathrm{Ext}_{\mathrm{A}}^{i}(\mathrm{M},\mathrm{N})\) 的支撑集亦然。
\item
  设 A 是 Noether 的且模 M 与 N 都有限生成；若 A-模 M⊗A N 具有有限长度，则 Tor \(_{i}^{A}\) (M, N) 与 Ext \(_{A}^{i}\) (M, N) 亦然。
\end{enumerate}

若 p 是 A 的一个不属于 Supp(M) ∩ Supp(N) 的素理想，则模 \(M_{p}\) 或 \(N_{p}\) 之一为零，于是由命题 2 得到 a)。

为使 Noether 环上的有限生成模具有有限长度，必须且只需其支撑集由极大理想组成 (IV, § 2, 第 5 节, 命题 7)。在假设 b) 下，A-模 Tor \(_{i}^{A}\) (M, N) 与 Ext \(_{A}^{i}\) (M,N) 都是有限生成的 (A, X, 第 108 页, 推论)；因此断言 b) 由 a) 得出。

命题 3.--- 设 A 为 Noether 环，M 为有限生成 A-模，N 为 A-模。

\begin{enumerate}
\def\labelenumi{\alph{enumi})}
\tightlist
\item
  我们有
\end{enumerate}

\[
\mathrm{dp} _ {\mathrm{A}} (\mathrm{M}) = \sup _ {\mathfrak {p}} \mathrm{dp} _ {\mathrm{A} _ {\mathfrak {p}}} \left(\mathrm{M} _ {\mathfrak {p}}\right), \quad \mathrm{di} _ {\mathrm{A}} (\mathrm{N}) = \sup _ {\mathfrak {p}} \mathrm{di} _ {\mathrm{A} _ {\mathfrak {p}}} \left(\mathrm{N} _ {\mathfrak {p}}\right), \quad \mathrm{dh} (\mathrm{A}) = \sup _ {\mathfrak {p}} \mathrm{dh} \left(\mathrm{A} _ {\mathfrak {p}}\right),
\]

其中 p 取遍 A 的素理想（相应地极大理想）组成的集合。

\begin{enumerate}
\def\labelenumi{\alph{enumi})}
\setcounter{enumi}{1}
\tightlist
\item
  映射 \(\mathfrak{p} \mapsto \mathrm{dp}_{A_{\mathfrak{p}}}(\mathrm{M}_{\mathfrak{p}})\)（从 \(\operatorname{Spec}(A)\) 到 \(\overline{Z}\)）是上半连续的。
\end{enumerate}

我们来证明 a)。设 \(n\) 为一个 \(\geqslant 0\) 的整数。设 \(\mathrm{dp}_{A}(\mathbf{M}) < n\)。对 A 的每个素理想 \(\mathfrak{p}\) 及每个 \(\mathrm{A}_{\mathfrak{p}}\) -模 Q，\(\mathrm{A}_{\mathfrak{p}}\) -模 \(\mathrm{Ext}_{\mathrm{A}_{\mathfrak{p}}}^{n}(\mathrm{M}_{\mathfrak{p}},\mathrm{Q})\) 同构于 \((\mathrm{Ext}_{\mathrm{A}}^{n}(\mathrm{M},\mathrm{Q}))_{\mathfrak{p}}\) (命题 2)，故为零；由此推出不等式 \(\mathrm{dp}_{A_{\mathfrak{p}}}(\mathrm{M}_{\mathfrak{p}}) \leqslant \mathrm{dp}_{\mathrm{A}}(\mathrm{M})\) (A, X, 第 134 页, 命题 1)。反之，设有 \(\mathrm{dp}_{A_{\mathfrak{m}}}(\mathrm{M}_{\mathfrak{m}}) < n\)（对 A 的每个极大理想 \(\mathfrak{m}\)），并设 R 为一个 A-模。有 \((\operatorname{Ext}_{\mathrm{A}}^{n}(\mathrm{M}, \mathrm{R}))_{\mathfrak{m}} = 0\)（对每个 \(\mathfrak{m}\)）(命题 2)，故 \(\operatorname{Ext}_{\mathrm{A}}^{n}(\mathrm{M}, \mathrm{R}) = 0\) (II, § 3, 第 3 节, 定理 1 的推论 2)，这蕴含 \(\mathrm{dp}_{\mathrm{A}}(\mathrm{M}) < n\) (A, X, 第 134 页, 命题 1)。a) 的第一个等式由此得出。第二个等式同样证明，利用命题 1 (iii) 给出的内射维数的刻画。由于 dh(A) 是诸 A-模的内射维数组成的集合的上确界 (第 1 节)，第三个等式由此得出。

我们来证明 b)。设 \(\mathfrak{p}\) 为 A 的一个素理想，\(n = \mathrm{dp}_{\mathrm{A}_{\mathfrak{p}}}(M_{\mathfrak{p}})\)。我们来证明存在 \(\mathfrak{p}\) 在 \(\operatorname{Spec}(A)\) 中的一个邻域 U，使得有 \(\mathrm{dp}_{\mathrm{A}_{\mathfrak{q}}}(M_{\mathfrak{q}}) \leqslant n\)（对每个 \(\mathfrak{q} \in U\)）。当 \(n = +\infty\) 时这是显然的；当 \(n = -\infty\) 时，这由 M 的支撑集是闭集这一事实得出。现设 \(n\) 有限，并选取 A-模的一个正合列

\[
\mathrm{P} _ {n - 1} \stackrel {d _ {n - 1}} {\longrightarrow} \mathrm{P} _ {n - 2} \longrightarrow \dots \longrightarrow \mathrm{P} _ {0} \stackrel {d _ {0}} {\longrightarrow} \mathrm{M} \rightarrow 0,
\]

其中诸 \(P_{i}\) 是有限型自由模 (A, X, 第 53 页, 命题 6)。令 \(P = Ker d_{n-1}\)；它是一个有限表示模。\(A_{p}\) -模 \(P_{p}\) 是投射的 (A, X, 第 134 页, 命题 1)，因而是自由的 (II, § 3, 第 2 节, 命题 5 的推论 2)。由 II, § 5, 第 1 节, 命题 2 的推论，存在 \(A - p\) 的一个元素 f 使得 \(A_{f}\) -模 \(P_{f}\) 是自由的；于是对 Spec(A) 中由不含 f 的素理想组成的开集 U 的每个素理想 q，\(A_{q}\) -模 \(P_{q}\) 都是自由的。这就证明了 b)。

推论 1.--- 对 A 的每个乘性子集 S，我们有

\[
\mathrm{dp} _ {\mathrm{S} ^ {- 1} \mathrm{A}} (\mathrm{S} ^ {- 1} \mathrm{M}) \leqslant \mathrm{dp} _ {\mathrm{A}} (\mathrm{M}), \quad \mathrm{di} _ {\mathrm{S} ^ {- 1} \mathrm{A}} (\mathrm{S} ^ {- 1} \mathrm{N}) \leqslant \mathrm{di} _ {\mathrm{A}} (\mathrm{N}), \quad \mathrm{dh} (\mathrm{S} ^ {- 1} \mathrm{A}) \leqslant \mathrm{dh} (\mathrm{A}).
\]

推论 2.--- 若对 Supp(M) 的每个极大理想 m 有 dp \(_{A_m}\) (M \(_m\) ) \textless{} +∞，则 dp \(_A\) (M) \textless{} +∞。

事实上，由极大理想组成的 Supp(M) 的子空间 X 是拟紧的 (II, § 4, 第 2 节, 命题 8 与 9 及第 3 节, 命题 11 的推论 7)；映射 \(\mathfrak{m} \mapsto \mathrm{dp}_{\mathrm{A}_{\mathfrak{m}}}(\mathrm{M}_{\mathfrak{m}})\)（从 X 到 \(\overline{\mathbf{R}}\)）是上半连续的 (命题 3)，故有界 (TG, IV, 第 30 页, 定理 3 的推论)。

注. --- 设 A 为一个无限维的正则 Noether 环 (VIII, § 5, 习题 6 c))。下文将看到 (第 7 节, 定理 2 及 § 4, 第 1 节, 命题 1)，对 A 的每个极大理想 m 有 \(\mathrm{di}_{\mathrm{A_m}}(\mathrm{A_m}) = \mathrm{dh}(\mathrm{A_m}) < +\infty\)；于是命题 3 蕴含 \(\mathrm{di}_{\mathrm{A}}(\mathrm{A}) = \mathrm{dh}(\mathrm{A}) = +\infty\)。因此，函数 \(\mathfrak{p} \mapsto \mathrm{di}_{A_{\mathfrak{p}}}(\mathrm{N}_{\mathfrak{p}})\) 与 \(\mathfrak{p} \mapsto \mathrm{dh}(\mathrm{A}_{\mathfrak{p}})\) 一般而言不是上半连续的。

